% !TeX spellcheck = en_US
% Chapter 2
\chapter{بررسی ادبیات و مبانی نظری}\label{chap:review}

\section{مفاهیم پایه}\label{chap:brain}

\subsection{مقدمه}
دامنه بررسی (از تئوری DT تا کاربرد LLM) و روش مرور (systematic review بر اساس مقالات IEEE, ScienceDirect, و arXiv تا ۲۰۲۵).

\subsection{دوقلوی دیجیتال}
تعاریف (Grieves, Tao et al.)، اجزا (فیزیکی، مجازی، ارتباطی)، و تکامل تا Industry ۵.۰.



\subsection{کاربردهای دوقلوی دیجیتال}
در هوافضا، تولید، و نگهداری پیش‌بینی‌کننده (با تمرکز بر موتورهای جت).

\subsection{عامل هوشمند}
عریف (autonomy, reactivity)، معماری‌ها (BDI).

\subsection{سیستم‌های چندعاملی}
مفاهیم همکاری، پلتفرم‌ها (JADE)، و نقش در سیستم‌های توزیع‌شده.

\subsection{دوقلوهای دیجیتال چندعاملی}
ادغام MAS با DT برای تصمیم‌گیری پویا (بر اساس surveyهای ۲۰۲۴ مانند "The Role of Multi-Agents in Digital Twin Implementation")

\subsection{هوش مصنوعی در دوقلوهای دیجیتال}
نقش AI در شبیه‌سازی و پیش‌بینی.

\subsection{یادگیری عمیق و کاربرد آن در دوقلوهای دیجیتال}
تکنیک‌های DL برای پردازش داده‌های زمانی.


\subsection{لایه توجه و سازوکار آن در مدل‌های یادگیری عمیق}
توضیح attention mechanism در transformerها و کاربرد در پیش‌بینی سری‌های زمانی.


\subsection{شبکه‌های عصبی بازگشتی و کاربرد آن در دوقلوهای دیجیتال}
معرفی RNN, LSTM, GRU، و استفاده در پیش‌بینی RUL (با مثال از CMAPSS).


\subsection{مدل‌های زبانی بزرگ}
تعریف \glspl{LLM}


تعریف LLM (مانند GPT-4o در ۲۰۲۵)، قابلیت‌های reasoning، و ابزارهای workflow مانند Dify یا LangChain برای ادغام با DT.

\subsection{جمع‌بندی}
ارتباط مفاهیم با معماری پیشنهادی شما.


\section{مرور کارهای مرتبط}

\subsection{مقالات مرتبط با دوقلوهای دیجیتال}

بررسی surveyهای کلی مانند "Digital twin: Data exploration, architecture, implementation and future" (۲۰۲۴) و "A new era for digital twins: progress and industry adoption" (۲۰۲۵). جدول مقایسه کاربردها.

\subsection{مقالات مرتبط با معماری دوقلوهای دیجیتال}

تمرکز بر معماری‌های لایه‌ای مانند "Smart City Digital Twins: A Modular and Adaptive Architecture" (۲۰۲۴) و "Architectural Design for Digital Twin Networks" (۲۰۲۵). تحلیل چگونگی الهام‌بخشی به معماری ۵ لایه‌ای شما.

\subsection{مقالات مرتبط با دوقلوهای دیجیتال چندعاملی}
کارهای اخیر مانند "Advancing Robotic Systems with Distributed Multi-Agent Digital Twins" (۲۰۲۵)، "Multi-Agent Systems for Manufacturing Digital Twins" (۲۰۲۵)، و "Agent-based digital twins for collaborative machine intelligence" (۲۰۲۵). بحث در مورد کاربرد در هوافضا (مانند "Developing a Digital Twin for Testing Multi-Agent Systems in Advanced Air Mobility"، ۲۰۲۵).

\subsection{مقالات مرتبط با کاربردهای مدل‌های زبانی بزرگ در دوقلوهای دیجیتال}

مقالات ۲۰۲۴-۲۰۲۵ مانند "Integrating large language model and digital twins in the context of Industry 5.0" (۲۰۲۵)، "Leveraging Large Language Models for Enhanced Digital Twin Modeling" (۲۰۲۵)، و "Large language models-enabled digital twins for precision medicine" (۲۰۲۵). تأکید بر ادغام برای تعامل هوشمند.


\subsection{مقالات مرتبط با داده ساختار \lr{CMAPSS}}
کارهای پیش‌بینی RUL مانند "Linear Methods for Predictive Maintenance: The Case of NASA C-MAPSS" (۲۰۲۵)، "A deep learning-based prognostic approach for predicting turbofan engine RUL" (۲۰۲۵)، و "Remaining useful life prediction of turbofan engines using LSTM" (۲۰۲۵). مقایسه با روش‌های DL و پیشنهاد بهبود با MAS و LLM.

\subsection{تحلیل شکاف‌ها و نتیجه‌گیری فصل}
شناسایی کمبودها (مانند عدم ادغام LLM در MAS-DT برای توضیح‌پذیری) و جایگاه نوآوری شما.





%برای بیان اکرونیم یک کلمه مانند  \gls{LASSO}، به این صورت عمل کنید.
%   \gls{ML}
%   یا 
%   \gls{DL}
%   یا
%   \gls{FN}
%   یا 
%   \gls{TN}
%   یا
%   \gls{LSTM}
%   یا 
%   \gls{SGL}
%   یا
%   \gls{TP}
%   یا 
%   \gls{FP}

%\subsection{زیربخش \gls*{1st}}\label{sec:brain:sub_brain}
%این یک متن نمونه است که به سه مرحله \gls{1st}،
% \gls{2nd} و \gls{3rd} تقسیم می‌شود. 
% 
% برای معادل جمع فارسی در دیکشنری اینطور فراخوان کنید: \glspl{2nd}
%\begin{itemize}
%	\item
%	\gls{1st}: 
%	این یک آیتم نمونه است...
%		\item 
%	\gls{2nd}: 
%	این یک آیتم نمونه است...
%	\item 
%	\gls{3rd}: 
%	این یک آیتم نمونه است...
%\end{itemize}
%
%\subsection{زیربخش دوم}\label{sec:brain:second_sub}
%این یک متن نمونه است...
%
%
%\paragraph{پاراگراف نمونه \gls*{1st}}
%این یک متن نمونه است...
%\paragraph{پاراگراف نمونه \gls*{2nd}}
%این یک متن نمونه است...
%
%\section{معرفی \gls*{dataset}}\label{sec:dataset}
%این بخش نمونه است...
%
%\begin{enumerate}
%	
%\item\textbf{فاز اول یا \lr{English-1}:} 
%این یک متن نمونه است.
%\begin{itemize}
%	\item 
%	این یک آیتم نمونه است...
%	\item 
%	این یک آیتم نمونه است...
%	\item 
%	این یک آیتم نمونه است...
%\end{itemize}
%
%این یک متن نمونه است...
%
%\item\textbf{فاز دوم یا \lr{English-2}:} 
%این یک متن نمونه است...
%
%\item\textbf{فاز سوم یا \lr{English-3}:} 
%این یک متن نمونه است...
%\end{enumerate}
%
%برای کامنت کردن بخشی از پروژه بدون حذف آن، از این دستور استفاده کنید.
%\begin{comment}
%برای کامنت کردن بخشی از پروژه بدون حذف آن، از این دستور استفاده کنید.
%\end{comment}
%
%برای درج جدول \ref{tab:my_table1} و رفرنس‌دهی از این کد استفاده کنید.
%
%\begin{table}[t!]
%	\centering
%	\caption{جدول نمونه افقی}
%	\label{tab:my_table1}
%%	\begin{tabular}{rccccc}
%%		\toprule			
%%		$\text{گروه}$ && $\text{فارسی (English)}$ & $\text{فارسی English}$ & $\text{فارسی English}$ & $\text{فارسی English}$ \\
%%		\midrule
%%		\rl{جنسیت}&(\rl{مرد - زن})	
%%		& $10 - 20$ & $10 - 20$ & $10 - 20$ & $10 - 20$
%%		\\
%%		\bottomrule
%%	\end{tabular}
%\end{table}
%
%برای درج جدول عمودی \ref{tab:my_table2} و رفرنس‌دهی از این کد استفاده کنید.
%\begin{sidewaystable}
%	\begin{threeparttable}
%		\caption{جدول نمونه عمودی}
%		\label{tab:my_table2}
%		\setlength\tabcolsep{0pt} % make LaTeX figure out intercolumn spacing
%		
%%		\begin{tabular*}{\columnwidth}{@{\extracolsep{\fill}} cccccccccc}
%%			\toprule			
%%			$\text{گروه}$ & 
%%			$\myfrac{\text{جنسیت}}{\text{مرد - زن}}$ & 
%%			$\myfrac{\text{فارسی}}{\text{\lr{mean$\pm $std}}}$ & 
%%			$\myfrac{\text{ٍenglish}}{\text{\lr{mean$\pm $std}}}$ & 
%%			$\myfrac{\text{فارسی}}{\text{\lr{mean$\pm $std}}}$ & 
%%			$\myfrac{\text{\lr{english}}}{\text{\lr{mean$\pm $std}}}$ & 
%%			$\myfrac{\text{\lr{ٍenglish}}}{\text{\lr{mean$\pm $std}}}$ & 
%%			$\myfrac{\text{ٍenglish}}{\text{\lr{mean$\pm $std}}}$ & 
%%			$\myfrac{\text{\lr{ٍenglish}}}{\text{\lr{mean$\pm $std}}}$ & 
%%			$\myfrac{\text{\lr{RAVALT-PF}}}{\text{\lr{mean$\pm $std}}}$ \\
%%			\midrule
%%			فارسی AD & $189 / 153$ & $75.0 \pm 7.8$ & $4.39 \pm 1.67$ & $15.2 \pm 3.0$ & $19.6 \pm 6.7$ & $29.9 \pm 8.1$ & $23.2 \pm 2$ & $22.8 \pm 7.6$ & $89.0 \pm 21.2$ 
%%			\\
%%			\bottomrule
%%		\end{tabular*}
%	\end{threeparttable}
%\end{sidewaystable}
%
%\section{فرمول}\label{sec:formula}
%بنابراین برای هر نقطه $X$ از تصویر سه بعدی دریافتی، رابطه
%\begin{align} \label{eq:n3-main}
%	v(X) = u(X) f(X) + n(X)
%\end{align}
%برقرار است که $v$ سیگنال اندازه‌گیری شده، $u$ سیگنال مورد انتظار بدون اثر \gls{bias} میدان و $f$ یک سیگنال با تغییرات آرام و نماینده \gls{bias} میدان است.
%روابط تبدیل فوریه تابع $f(t)$ و معکوس آن به بیان شده در روابط \ref{eq:fourier} و \ref{eq:afourier}  هستند.
%\begin{align}\label{eq:fourier}
%\mathcal{F}(f(t))= F(w)=\sum_{k=0}^{N-1}f_k e^{-i\frac{2\pi}{N}kn}
%\end{align}
%\begin{align}\label{eq:afourier}
%f_k= F(w)=\frac{1}{N}\sum_{n=0}^{N-1}\mathcal{F}_n e^{-i\frac{2\pi}{N}kn}
%\end{align}

